% Standalone source-record documentation; no research attempts.
\documentclass[11pt]{article}
\usepackage[margin=25mm]{geometry}
\usepackage{fontspec,unicode-math}
\setmainfont{Latin Modern Roman}
\setmathfont{Latin Modern Math}
\usepackage{amsmath,amssymb,mathtools,xurl,hyperref}
\hypersetup{colorlinks=true,urlcolor=blue,linkcolor=blue}
\setcounter{secnumdepth}{0}
\setlength{\parindent}{0pt}
\setlength{\parskip}{5pt}
\setlength{\emergencystretch}{3em}
\newcommand{\sref}[2]{\hyperlink{source:#1:#2}{[S#2]}}
\newcommand{\cataloguescope}[1]{\paragraph{Source record.}#1}
\newcommand{\floor}[1]{\left\lfloor#1\right\rfloor}
\begin{document}
% UnsolvedMath ID 30000782; source catalogue code OWR-1587-001
\setcounter{section}{30000781}
\section{Boolean Evaluation and Arithmetic-Circuit Complexity}\label{30000782}
{\small\sffamily UnsolvedMath ID 30000782 · Computer Science}
\subsection{Definitions and mathematical statement}

\paragraph{Shared notation.}Standard number systems are $\mathbb N,\mathbb Z,\mathbb Q,\mathbb R,\mathbb C$. Specialized notation and index ranges follow the cited source and the definitions below.


Polynomial time means a number of elementary steps bounded by a polynomial in the encoded input length. A decision problem is in $\mathrm{NP}$ if yes-instances have polynomial-length certificates verifiable in polynomial time; NP-hardness is under polynomial-time reductions.

An arithmetic circuit is a directed acyclic graph whose input nodes are variables or field constants and whose internal nodes perform additions and multiplications (with subtraction allowed when stated). Size counts gates; depth is the maximum number of gates on an input-output path.
\paragraph{Record scope and qualifications.}
The catalogue title is retained as metadata. Named certificates, contracts, bounds or frameworks in that title are not asserted as results of the original source extract.
The supplied statement refers to one polynomial without specifying a growing polynomial family, degree bounds, or coefficient encoding. The intended asymptotic comparison is that of the source.

Source-record limitation: The required growing polynomial family, degree/coefficient encoding, and complexity-size parameters are absent; a single fixed polynomial does not define the stated asymptotic target.
\paragraph{Statement recorded in the supplied source.}
If a multivariate polynomial $P$ can be evaluated on rational inputs by a Boolean polynomial-time algorithm, must $P$ have a polynomial-size arithmetic circuit using only additions, subtractions, and multiplications? \sref{30000782}{2}
\paragraph{Source status.}
The supplied catalogue labels this source record open; that metadata was checked against the supplied snapshot on 2026-10-11. The label is preserved for provenance and is not an independent verification of present open status or mathematical correctness. The catalogue literature assessment is dated 2026-08-21.
\subsection{Short English statement}
Determine whether efficient Boolean evaluation of a polynomial entails a small arithmetic circuit.
\subsection{Sources}
\begin{itemize}
\item[\hypertarget{source:30000782:1}{S1}] Ulam AI and the UnsolvedMath Contributors, UnsolvedMath: A Curated Collection of Open Mathematics Problems, record 30000782 (OWR-1587-001). Catalogue curation and metadata: CC BY 4.0. The source extract is retained; definitions, English summaries and qualifications are editorial additions.
\url{https://huggingface.co/datasets/ulamai/UnsolvedMath}.
\item[\hypertarget{source:30000782:2}{S2}] Complexity Theory (2008), Oberwolfach Reports, DOI 10.4171/owr/2007/31, pp. 1793--1864
\url{https://doi.org/10.4171/owr/2007/31}.
\item Source-specific Oberwolfach Reports licences, including CC BY-SA 4.0 where applicable, are retained; see THIRD\_PARTY\_NOTICES.md.
\end{itemize}
\paragraph{Preparation and review.}
The original source and UnsolvedMath contributors retain attribution for the source extract. Codex (GPT-6; exact serving revision not exposed) assisted with definitions, English summaries, scope notes and formatting. No new human mathematical authorship or independently verified solution is asserted. Permanent source ID 30000782 and code OWR-1587-001 checked against the complete supplied registry and current repository. Distinct upstream records, including repeated questions and broad or damaged extracts, are retained for complete coverage of the source CS category. A repeated record is not a new independent conjecture. Complete mathematical-equivalence and current-literature review has not been performed. Preparation and final review follow contribution prompts C and D, with source gaps disclosed. Structural and rendering checks do not establish a new theorem or open status.
\label{end:30000782}
\end{document}
